\documentclass[11pt]{article}
\usepackage[margin=1in]{geometry}
\usepackage{amsmath,amssymb}
\usepackage[hidelinks]{hyperref}

\begin{document}

\section*{Human Review: Selector-Block Dual-Span Counterexample}

\noindent Original automatic proof: \texttt{solution\_packet.pdf}

\noindent Checked date: 19 June 2026

\noindent Status: Manually checked.

\noindent Verdict: The proof appears correct.

\section*{Human Comments}

The packet gives a counterexample to Question 6.1 of Borodulin-Nadzieja, Farkas, Jachimek, and Pelczar-Barwacz, concerning whether the norm-closed span, or even the norm-closed convex hull, of the sign functionals
\[
  W(\overline{\mathcal F})
\]
recovers the full dual space or the full dual unit ball of the corresponding combinatorial Banach space.

The construction is clear and appears sound.  One decomposes the underlying set into finite blocks
\[
  \Omega=\bigsqcup_{k=1}^\infty I_k,
  \qquad |I_k|=m_k\to\infty,
\]
and takes \(\mathcal F\) to be the hereditary family of finite selectors, namely finite sets meeting each block in at most one point.  Its closure consists exactly of the possibly infinite selectors
\[
  A\subseteq\Omega,
  \qquad |A\cap I_k|\leq 1\quad(k\in\mathbb N),
\]
so every element of \(W(\overline{\mathcal F})\) is a signed selector with at most one nonzero coordinate in each block.

The identification of the associated space is correct:
\[
  \|x\|_{\mathcal F}
  =\sum_k \|x|_{I_k}\|_\infty,
\]
and hence
\[
  X_{\mathcal F}
  =\left(\bigoplus_k \ell_\infty^{m_k}\right)_{\ell_1},
  \qquad
  X_{\mathcal F}^*
  =\left(\bigoplus_k \ell_1^{m_k}\right)_{\ell_\infty}.
\]
The dual norm is therefore
\[
  \|\alpha\|=\sup_k \|\alpha|_{I_k}\|_1.
\]

The obstruction is the blockwise uniform vector \(u\in X_{\mathcal F}^*\), defined by
\[
  u(\omega)=\frac1{m_k}\qquad(\omega\in I_k).
\]
For every block, \(\|u|_{I_k}\|_1=1\), so \(u\) belongs to the dual unit ball.

On the other hand, if
\[
  y=\sum_{j=1}^r a_j\sigma_j,
  \qquad \sigma_j\in W(\overline{\mathcal F}),
\]
then in each block \(I_k\), the vector \(y|_{I_k}\) is supported on at most \(r\) coordinates, because each signed selector contributes to at most one coordinate in that block.  Choosing \(k\) with \(m_k>2r\), one obtains
\[
  \|u|_{I_k}-y|_{I_k}\|_1
  \geq \frac{m_k-r}{m_k}
  >\frac12.
\]
Therefore
\[
  \|u-y\|_{X_{\mathcal F}^*}>\frac12
\]
for every finite linear combination \(y\) of elements of \(W(\overline{\mathcal F})\).  This proves that \(u\) is not in the norm-closed linear span of \(W(\overline{\mathcal F})\).  Since the convex hull is contained in the linear span, the same vector also witnesses failure of the asserted norm-closed convex-hull equality.

The proof is short, but the mechanism is convincing: finite combinations of selectors are uniformly sparse in each block, whereas the dual unit ball contains vectors whose mass is spread over arbitrarily large blocks.  The growth \(m_k\to\infty\) is exactly what makes the obstruction uniform over all finite combinations.

\section*{Review Outcome}

I would classify this as a correct counterexample to the general form of Question 6.1.  It should be noted that the construction does not address possible positive results for more restrictive subclasses of hereditary families, such as compact families or other special combinatorial families.  For the unrestricted formulation considered in the packet, however, the proof appears correct.

\end{document}
